\chapter{Utilisation de l'intelligence artificielle (IA)}
\label{ann:ia}

Le tableau~\ref{tab:usage_ia} détaille les modèles employés et l'usage qui en a été fait. Le périmètre général et les vérifications effectuées sont exposés dans la déclaration figurant en tête de ce document.

\begin{table}[H]
    \centering
    \renewcommand{\arraystretch}{1.4}
    \caption{Utilisation des modèles d'IA}
    \label{tab:usage_ia}
    \begin{tabular}{p{0.22\linewidth} p{0.56\linewidth} >{\centering\arraybackslash}p{0.13\linewidth}}
        \hline
        \textbf{Modèle utilisé}                 & \textbf{Utilisé pour ...}                                                                                                    & \textbf{Chapitres concernés} \\
        \hline
        Claude Opus 4.8 \newline + Infomaniak Euria                        & Aide à la rédaction, à la relecture critique et à la reformulation du texte, correction de l'orthographe et de la grammaire. & Tous les chapitres           \\
        Claude Opus 4.8 \newline + Infomaniak Euria                        & Aide au raccourcissement sections trop longues & Tous les chapitres           \\
        Claude Opus 4.8                         & Aide à la création de diagrammes et de schémas DRAWIO & 3, 4, 5           \\
        Claude Opus 4.8                         & Aide à la recherche de sources et de références, et à la synthèse des informations recueillies.                              & 2                            \\
        Claude Opus 4.8                         & Aide à la décomposition du système et à la séparation des blocs. & 3                            \\
        Claude Opus 4.8 \newline + Gemini Pro 3 & Aide à la recherche de composants, à la lecture des datasheets et à la vérification de compatibilité entre composants.       & 4                            \\
        Claude Opus 4.8                         & Aide à la relecture du schéma électrique et du layout.                                                                       & 4                            \\
        OpenAI GPT-5.3-Codex                    & Écriture du code de tracé des courbes de consommation.                                                                       & Annexe~\ref{ann:conso_esp32} \\
        Claude Opus 4.8                         & Aide au choix de l'architecture logicielle.                                                                                  & 5                            \\
        Claude Opus 4.8                         & Aide à l'écriture, à la relecture et à la revue du firmware, et à la mise en place des tests automatisés.                    & 5                            \\
        Infomaniak Euria                         & Aide à la création de scripts Python pour générer les plots                & Annexes \ref{ann:endurance}, \ref{ann:mesures} et \ref{ann:autonomie}                            \\
        \hline
    \end{tabular}
\end{table}